\clearpage
\section{The Weyl algebra}\label{sec:weyl}

\noindent\motiv{\vet{The observables of the free field are generated by the Weyl operators
$W(u)$, formally $e^{\ii\phi(u)}$, labelled by the vectors $u$ of a real vector space $V$ with an
alternating form $\sigma$, the classical phase space of the field. On the Hilbert space
$\ell^2(V)$ these operators are unitary, and the C*-algebra they generate serves as the algebra
$A$ of the axioms for the free field.} \revc{The Hilbert space $\ell^2(V)$ carries the regular
representation of the additive group $V$, twisted by the phases of $\sigma$. It serves to show that
Weyl operators exist; it is not the vacuum (Fock) representation of the field, which these notes do
not construct (Remarks~\ref{rem:separating} \revd{and~\ref{rem:fock}}).}}

\begin{remark}[Signs]\label{rem:dictionary}
\vet{Spacetime signature is mostly plus, and the Klein--Gordon operator is $P=-\square+m^2$
(\rev{$=\partial_t^2-\nabla^2+m^2$} in flat space). The \revd{causal propagator} is
$E=E^{\mathrm{adv}}-E^{\mathrm{ret}}$ for this $P$, so that
\rev{$[\phi(f),\phi(g)]=\ii\langle f,Eg\rangle$} and, at equal times, $[\phi(t,x),\Pi(t,y)]=\ii\delta(x-y)$ for the momentum $\Pi=\partial_t\phi$.\extraonly{ \rev{On the lattice this sign is
Proposition~\ref{prop:commutator}(e).}}
With $[f]$ the class of a test function $f$ modulo the kernel of $E$,
\rev{$\sigma([f],[g])=\langle f,Eg\rangle$} and $W(u)=e^{\ii\phi(u)}$, the Baker--Campbell--Hausdorff formula
$e^{X}e^{Y}=e^{X+Y}e^{[X,Y]/2}$ for a central commutator gives
$W(u)W(w)=e^{-\ii\sigma(u,w)/2}W(u+w)$, which is Proposition~\ref{prop:weyl-relation} below.}
\end{remark}

\subsection{Weyl operators on \texorpdfstring{$\ell^2(V)$}{\textell\texttwosuperior(V)}}

\noindent Throughout, $V$ is a real vector space of any dimension, which we regard as a discrete
set, and $\sigma$ is a real bilinear form on $V$ that is \emph{alternating}: $\sigma(u,u)=0$,
hence $\sigma(w,u)=-\sigma(u,w)$. $\sigma$ may be degenerate. $\ell^2(V)$ is the Hilbert space
of square-summable $\psi:V\to\C$, and $\delta_v$ is the indicator function of $\{v\}$.

\begin{definition}[Weyl operator]\label{def:weyl}
For $u\in V$, $(W(u)\psi)(v)=e^{\ii\sigma(u,v)/2}\,\psi(v+u)$.
\end{definition}

\begin{lemma}[Isometry]\label{lem:weyl-isometry}
\revk{The Weyl operators on $\ell^2(V)$ (Definition~\ref{def:weyl}) satisfy} $\|W(u)\psi\|=\|\psi\|$
for \revk{all $u\in V$,} $\psi\in\ell^2(V)$, and $W(0)=1$.
\end{lemma}
\begin{proof}
$|(W(u)\psi)(v)|^2=|\psi(v+u)|^2$, and $v\mapsto v+u$ is a bijection of $V$, so the two sums of
squares have the same terms.
\end{proof}

\begin{proposition}[Weyl relation]\label{prop:weyl-relation}
\revk{For a real vector space $V$ with alternating form $\sigma$, the Weyl operators on $\ell^2(V)$
satisfy} $W(u)W(w)=e^{-\ii\sigma(u,w)/2}\,W(u+w)$ \revk{for all $u,w\in V$}.
\end{proposition}
\begin{proof}
$(W(u)W(w)\psi)(v)=e^{\ii\sigma(u,v)/2}e^{\ii\sigma(w,v+u)/2}\psi(v+u+w)$ and
$(W(u+w)\psi)(v)=e^{\ii\sigma(u+w,v)/2}\psi(v+u+w)$; by bilinearity the exponents differ by
$\ii\sigma(w,u)/2=-\ii\sigma(u,w)/2$.
\end{proof}

\begin{proposition}[Unitarity]\label{prop:weyl-unitary}
\revk{For $u\in V$ the Weyl operator $W(u)$ on $\ell^2(V)$ satisfies} $W(u)W(-u)=1$, $W(u)^*=W(-u)$, and $W(u)$ is unitary.
\end{proposition}
\begin{proof}
Proposition~\ref{prop:weyl-relation} and $\sigma(u,u)=0$ give $W(u)W(-u)=W(0)=1$. An isometry
$a$ satisfies $a^*a=1$, by polarization of $\langle a\psi,a\psi\rangle=\langle\psi,\psi\rangle$.
Hence $W(u)^*=W(u)^*W(u)W(-u)=W(-u)$, and $W(u)W(u)^*=W(u)W(-u)=1$.
\end{proof}

\begin{definition}[CCR algebra]\label{def:ccr}
$\CCR(V,\sigma)=\calg\{W(u):u\in V\}$, the C*-subalgebra of the bounded operators on
$\ell^2(V)$ generated by the Weyl operators.
\end{definition}

\begin{example}[One degree of freedom]\label{ex:one-dof}
For $V=\R^2$ with $\sigma((a,b),(a',b'))=ab'-ba'$, the operators \revd{$W(a,0)$ and $W(0,b)$
satisfy $W(a,0)W(0,b)=e^{-\ii ab}W(0,b)W(a,0)$} by Proposition~\ref{prop:weyl-relation} applied twice.
This is the Weyl form $e^{\ii a\hat q}e^{\ii b\hat p}=e^{-\ii ab}e^{\ii b\hat p}e^{\ii a\hat q}$
of $[\hat q,\hat p]=\ii$. On $\ell^2(\R^2)$, however, \rev{$\lambda\mapsto W(\lambda u)$} is not continuous: for
\rev{$\lambda\neq0$} and $u\neq0$, \rev{$W(\lambda u)\delta_0=\delta_{-\lambda u}$} is orthogonal to $\delta_0$, so
\rev{$\|W(\lambda u)\delta_0-\delta_0\|=\sqrt2$}. The unitary group \rev{$\lambda\mapsto W(\lambda u)$} is therefore not
strongly continuous, \revd{so it has no self-adjoint generator, since $\lambda\mapsto e^{\ii\lambda A}$ is
strongly continuous for every self-adjoint $A$~\cite[Theorem~VIII.7]{reed-simon}}: no field operator
$\phi(u)$ exists in this representation.
\end{example}

\begin{example}[Zero form]\label{ex:sigma-zero}
For $\sigma=0$ the Weyl operators commute, $W(u)W(w)=W(u+w)$, and $\CCR(V,0)$ is commutative:
the closure of the commutative algebra spanned by the $W(u)$ is commutative.
\end{example}

\subsection{Weyl systems}

\begin{definition}[Weyl system]\label{def:weyl-system}
A \emph{Weyl system} over $(V,\sigma)$ is a unital C*-algebra $A$ with a map $W:V\to A$ such
that each $W(u)$ is unitary, $W(u)W(w)=e^{-\ii\sigma(u,w)/2}W(u+w)$, and
$\calg\{W(u):u\in V\}=A$.
\end{definition}

\begin{lemma}[Consequences of the relation]\label{lem:weyl-system}
In a Weyl system $W(0)=1$, $W(u)^*=W(-u)$ and $W(u)W(w)=e^{-\ii\sigma(u,w)}W(w)W(u)$; in
particular $W(u)$ and $W(w)$ commute when $\sigma(u,w)=0$.
\end{lemma}
\begin{proof}
The relation with $u=w=0$ gives $W(0)^2=W(0)$, and multiplying by $W(0)^*$ gives $W(0)=1$.
Then $W(u)W(-u)=e^{\ii\sigma(u,u)/2}W(0)=1$, so $W(-u)=W(u)^{-1}=W(u)^*$. Writing both
$W(u)W(w)$ and $W(w)W(u)$ as multiples of $W(u+w)$ gives the phase
$e^{-\ii\sigma(u,w)/2}/e^{-\ii\sigma(w,u)/2}=e^{-\ii\sigma(u,w)}$.
\end{proof}

\begin{proposition}[The concrete system]\label{prop:ccr-system}
$\CCR(V,\sigma)$ with the Weyl operators is a Weyl system.
\end{proposition}
\begin{proof}
Unitarity and the relation are Propositions~\ref{prop:weyl-unitary}
and~\ref{prop:weyl-relation}. For generation, the inclusion of $\CCR(V,\sigma)$ into the bounded
operators is an isometric $*$-homomorphism with closed image; it maps the $*$-subalgebra of
$\CCR(V,\sigma)$ generated by the $W(u)$ onto the one of the bounded operators generated by
them, so the closure of the former is the preimage of $\CCR(V,\sigma)$, which is all of it.
\end{proof}

\begin{remark}[\revm{Uniqueness and covariance}]\label{rem:weyl-covariance}
\vet{The zero algebra satisfies the three conditions of Definition~\ref{def:weyl-system};
$\CCR(V,\sigma)$ is not the zero algebra, since $\ell^2(V)\neq0$. Covariance is not part of the definition: for
$\CCR(V,\sigma)$ it is Proposition~\ref{prop:relabel}, and for another Weyl system we have to assume a family of
$*$-isomorphisms with $\alpha_L(W(u))=W'(Lu)$. \revm{For nondegenerate $\sigma$
(Definition~\ref{def:nondegenerate}) and both algebras nonzero such a family exists by the uniqueness
theorem for the CCR algebra~\cite{slawny} (see also~\cite{manuceau}
and~\cite[Thm.~5.2.8]{bratteli-robinson}), which~\cite[p.~225]{dimock} uses by citation; these
notes do not prove it, and it is not formalized. For degenerate $\sigma$ it can fail:} for $\sigma=0$ and $V\neq0$, $A=\C$ with $W\equiv1$ is a Weyl system, and
$\CCR(V,0)$ is not isomorphic to it, since the vectors $W(u)\delta_0=\delta_{-u}$
(Proposition~\ref{prop:right-weyl}) show that $\CCR(V,0)$ is not one-dimensional.}
\end{remark}

\subsection{Relabelling}

\noindent Let $(V',\sigma')$ be a second space and $L:V\to V'$ a linear bijection that is
\emph{symplectic}: $\sigma'(Lu,Lv)=\sigma(u,v)$.

\begin{definition}[Relabelling unitary]\label{def:relabel}
$(U_L\psi)(v')=\psi(L^{-1}v')$ for $\psi\in\ell^2(V)$ and $v'\in V'$.
\end{definition}

\begin{proposition}[Relabelling]\label{prop:relabel}
\revk{Let $L:V\to V'$ be a linear bijection with $\sigma'(Lu,Lv)=\sigma(u,v)$ and $U_L$ its
relabelling unitary (Definition~\ref{def:relabel}). Then:}
\begin{enumerate}[label=(\alph*),leftmargin=2.5em]
\item $U_L:\ell^2(V)\to\ell^2(V')$ is unitary, $U_L^*=U_{L^{-1}}$, $U_{\mathrm{id}}=1$ and
$U_{L_2L_1}=U_{L_2}U_{L_1}$.
\item $U_LW(u)U_L^*=W'(Lu)$, where $W'$ are the Weyl operators of $(V',\sigma')$.
\item $\alpha_L(a)=U_LaU_L^*$ restricts to a $*$-isomorphism
$\alpha_L:\CCR(V,\sigma)\to\CCR(V',\sigma')$ with $\alpha_L(W(u))=W'(Lu)$,
$\alpha_{\mathrm{id}}=\mathrm{id}$ and $\alpha_{L_2L_1}=\alpha_{L_2}\alpha_{L_1}$.
\end{enumerate}
\end{proposition}
\begin{proof}
\begin{enumerate}[leftmargin=2.5em]
\item[(a)] $U_L$ reindexes along the bijection $L^{-1}$, so it preserves norms, and $U_{L^{-1}}$ is its
inverse; an invertible isometry is unitary. The composition law is
$\psi(L_1^{-1}L_2^{-1}v'')=\psi((L_2L_1)^{-1}v'')$.
\item[(b)] With $v=L^{-1}v'$,
$(U_LW(u)U_L^*\psi)(v')=e^{\ii\sigma(u,v)/2}\psi(L(v+u))=e^{\ii\sigma'(Lu,v')/2}\psi(v'+Lu)$.
\item[(c)] Conjugation by $U_L$ is a continuous $*$-isomorphism of the bounded operators. By (b) it
maps the Weyl operators of $V$ onto those of $V'$, hence the generated $*$-algebra into
$\CCR(V',\sigma')$, and by continuity its closure as well; the same holds for $L^{-1}$, which
gives the inverse. The identities follow from (a).
\end{enumerate}
\end{proof}

\begin{remark}[Time reversal]\label{rem:antisymplectic}
\vet{If $\sigma'(Lu,Lv)=-\sigma(u,v)$, as for time reversal, no $*$-homomorphism \revd{$\eta$} with
\revd{$\eta(W(u))=W'(Lu)$} \rev{into a nonzero algebra} exists unless $\sigma=0$: applying \revd{$\eta$} to
Proposition~\ref{prop:weyl-relation} gives $e^{-\ii\sigma(u,w)/2}=e^{\ii\sigma(u,w)/2}$, while
\rev{$\sigma(\lambda u,w)=\lambda\sigma(u,w)$} takes every real value when $\sigma(u,w)\neq0$. This is the
algebraic form of Remark~\ref{rem:time-orientation}.} \rev{The relation does not exclude an antilinear
map, which conjugates the phase, as in Wigner's treatment of time reversal\revb{~\cite[Ch.~26]{wigner}}; such a map is not a
$*$-homomorphism of complex algebras and lies outside Definition~\ref{def:covariance}.}
\end{remark}

\subsection{Right Weyl operators and the vector state}

\noindent\motiv{\rev{The operators and the state of this section are tools for \revf{Chapter~\ref{sec:simple}}. The right Weyl operators commute with the whole algebra and move
$\delta_0$ to every basis vector, which makes the vector state at $\delta_0$ faithful
(Remark~\ref{rem:separating}); and $\tau(W(v))=[v=0]$ picks out the coefficient of $1$ in a
finite linear combination of Weyl operators.}}

\begin{definition}[Right Weyl operator]\label{def:right-weyl}
$(R(w)\psi)(v)=e^{-\ii\sigma(w,v)/2}\,\psi(v+w)$, the Weyl operator of the form $-\sigma$.
\end{definition}

\begin{proposition}[Right operators]\label{prop:right-weyl}
\revk{For $w\in V$ the right Weyl operator $R(w)$ on $\ell^2(V)$ (Definition~\ref{def:right-weyl})}
is unitary, $R(w)^*=R(-w)$, $R(w)R(w')=e^{\ii\sigma(w,w')/2}R(w+w')$, and $R(w)$
commutes with every element of $\CCR(V,\sigma)$. Moreover $W(u)\delta_0=\delta_{-u}$ and
$R(w)\delta_0=\delta_{-w}$.
\end{proposition}
\begin{proof}
The first three claims are Propositions~\ref{prop:weyl-relation} and~\ref{prop:weyl-unitary}
for $-\sigma$. The exponents of $(W(u)R(w)\psi)(v)$ and $(R(w)W(u)\psi)(v)$ differ by
$-\ii(\sigma(w,u)+\sigma(u,w))/2=0$, so $R(w)$ commutes with every $W(u)$, and so does
$R(w)^*=R(-w)$. By Lemma~\ref{lem:commute}, with $s$ the set of Weyl operators and
$s'=\{R(w)\}$, $R(w)$ commutes with $\CCR(V,\sigma)$. Finally
$(W(u)\delta_0)(v)=e^{\ii\sigma(u,v)/2}[v+u=0]$, and at $v=-u$ the phase is
$e^{-\ii\sigma(u,u)/2}=1$; the same for $R(w)$.
\end{proof}

\begin{definition}[Vector state]\label{def:vector-state}
$\tau(a)=\langle\delta_0,a\,\delta_0\rangle$ for $a\in\CCR(V,\sigma)$.
\end{definition}

\begin{proposition}[Values of the vector state]\label{prop:vector-state}
\revk{The vector state $\tau$ of $\CCR(V,\sigma)$ (Definition~\ref{def:vector-state})} is linear, $\tau(1)=1$, $\tau(a^*a)=\|a\delta_0\|^2$, $|\tau(a)|\leq\|a\|$, and
$\tau(W(v))=1$ if $v=0$ and $0$ otherwise.
\end{proposition}
\begin{proof}
$\tau(a^*a)=\langle a\delta_0,a\delta_0\rangle$, and $|\tau(a)|\leq\|a\delta_0\|\leq\|a\|$ by
the Cauchy--Schwarz inequality. By Proposition~\ref{prop:right-weyl},
$\tau(W(v))=\langle\delta_0,\delta_{-v}\rangle$.
\end{proof}

\begin{remark}[The tracial state]\label{rem:tracial}
\revd{$\tau$ is \revf{a} \emph{tracial state} of the CCR algebra~\cite{petz}:
$\tau(W(u)W(w))=e^{-\ii\sigma(u,w)/2}[u+w=0]$ is symmetric in $u$ and $w$, since
$\sigma(u,-u)=0$, so $\tau(ab)=\tau(ba)$ by linearity and continuity. The vector $\delta_0$ is
cyclic, since $W(u)\delta_0=\delta_{-u}$, so $\ell^2(V)$ is the GNS representation of $\tau$. For
any dynamics the KMS condition at $\beta=0$ reduces to $\omega(ab)=\omega(ba)$, so tracial
states are the equilibrium states at infinite temperature~\cite[Remark~2]{derezinski-pillet};
$\tau$ is not a vacuum.} \revf{For nondegenerate $\sigma$ it is the only tracial state: a tracial
state $\omega$ has $\omega(W(v))=\omega(W(u)W(v)W(u)^*)=e^{-\ii\sigma(u,v)}\omega(W(v))$}
\revg{by Lemma~\ref{lem:weyl-system}; for $v\neq0$ nondegeneracy gives $w$ with $\sigma(w,v)\neq0$,
and a real multiple $u$ of $w$ has $\sigma(u,v)\notin2\pi\Z$, so $\omega(W(v))=0=\tau(W(v))$. By
linearity and continuity $\omega=\tau$.}
\end{remark}

\begin{remark}[Faithfulness and reducibility]\label{rem:separating}
\rev{If $a\in\CCR(V,\sigma)$ and $a\delta_0=0$, then $a=0$: for every $w$,
$a\delta_{-w}=aR(w)\delta_0=R(w)a\delta_0=0$ by Proposition~\ref{prop:right-weyl}, and a bounded
operator vanishing on the orthonormal basis $\{\delta_v\}$ is zero. Hence $\tau(a^*a)=\|a\delta_0\|^2$
vanishes only for $a=0$. For $V\neq0$ the representation on $\ell^2(V)$ is reducible: for
$w\neq0$, $R(w)$ commutes with $\CCR(V,\sigma)$ but is not a multiple of $1$, since
$R(w)\delta_0=\delta_{-w}$ is orthogonal to $\delta_0$. Primitivity therefore needs another
representation\revf{ (Remark~\ref{rem:irrep}); Lemma~\ref{lem:faithful} restates
the first claim}.}
\end{remark}

\begin{remark}[Fock representations]\label{rem:fock}
\revd{A Fock representation is the usual alternative to $\ell^2(V)$. It needs a complex
structure on $V$ compatible with $\sigma$, that is, a choice of vacuum, which a curved
spacetime does not single out~\cite[p.~219]{dimock}. For nondegenerate $\sigma$ its Weyl
operators generate a C*-algebra isomorphic to $\CCR(V,\sigma)$ by the uniqueness theorem of
Remark~\ref{rem:weyl-covariance}, so the net and its axioms are the same. It differs in two
ways. It is irreducible, and $\lambda\mapsto W(\lambda u)$ is strongly continuous, so field
operators exist by Stone's theorem. And a bounded symplectic map is implemented by a unitary
only when its antilinear part is Hilbert--Schmidt, by Shale's theorem~\cite{shale}, whereas on
$\ell^2(V)$ every symplectic bijection is (Proposition~\ref{prop:relabel}). \revf{Theorem~\ref{thm:primitive}
needs an irreducible representation, and a Fock representation is one
(Remark~\ref{rem:irrep}).}}
\end{remark}
